\section{同じ選択を繰り返す}\label{節:部分積分の反復}

整式と指数関数，あるいは整式と三角関数の積では，一度の部分積分で整式の次数が一つ下がる。残る積分に同じ方針を適用すれば，やがて整式が定数になる。最初から一度で計算を終わらせようとする必要はない。何回か繰り返せば，確実に既知の積分へ到達することが見えている。

\begin{example}{二次式の次数を一つずつ下げる}
$\int x^2\cos x\dd x$ を求めよ。
\end{example}

まず $u=x^2$，$v'=\cos x$ と選ぶ。$u'=2x$，$v=\sin x$ なので，
\[
 \int x^2\cos x\dd x
 =x^2\sin x-2\int x\sin x\dd x.
\]
残った積分では，$x$ を微分し，$\sin x$ を積分する。$\sin x$ の原始関数は $-\cos x$ だから，
\[
\begin{aligned}
 \int x\sin x\dd x
 &=x(-\cos x)-\int1\cdot(-\cos x)\dd x\\
 &=-x\cos x+\sin x+C.
\end{aligned}
\]
これをもとの式へ戻して，任意定数をまとめると
\[
\begin{aligned}
 \int x^2\cos x\dd x
 &=x^2\sin x-2(-x\cos x+\sin x)+C\\
 &=x^2\sin x+2x\cos x-2\sin x+C.
\end{aligned}
\]
二回目の積分で現れた負号だけでなく，その式全体に最初の $-2$ が掛かることも確かめる。途中の積分を独立に計算したときは，戻す先の係数を残しておくと，符号の取り違えを防げる。

答えを微分すると，途中で現れた項が対になって消える。
\[
\begin{aligned}
 &\frac{\dd}{\dd x}
  \{x^2\sin x+2x\cos x-2\sin x\}\\
 &\qquad=2x\sin x+x^2\cos x
       +2\cos x-2x\sin x-2\cos x\\
 &\qquad=x^2\cos x.
\end{aligned}
\]
部分積分で補った項は，このように一つ前の微分で生まれた余分な項を打ち消している。繰り返す計算は，思い付きで項を増やす操作ではなく，積の微分による余分な項を順に取り除く操作である。

一般に，$P(x)$ が $n$ 次の整式なら，$P,P',P'',\ldots$ と微分を繰り返すと，$n$ 回目に定数になり，その次には $0$ になる。したがって $P(x)\e^{ax}$ や $P(x)\sin(ax)$，$P(x)\cos(ax)$ は，$a\ne0$ のとき，整式を微分する部分積分を繰り返して計算できる。$a$ が付く場合は，指数関数や三角関数を積分するたびに $a$ で割る。次数が下がることと，係数が変わることを別々に確認すればよい。

